\begin{figure}[H]
\centering
% Giới hạn theo chiều rộng A4; bố cục 2 cột nên chiều cao vừa một trang
\resizebox{0.98\textwidth}{!}{%
\begin{tikzpicture}[
    entity/.style={
        ellipse, draw=black, fill=white, line width=0.85pt,
        minimum width=1.9cm, minimum height=0.72cm,
        font=\scriptsize\bfseries, align=center, inner sep=0.8pt
    },
    process/.style={
        rectangle, draw=black, fill=gray!15, line width=0.95pt,
        minimum width=2.9cm, minimum height=0.9cm,
        font=\footnotesize\bfseries, align=center, rounded corners=2pt, inner sep=1.5pt
    },
    store/.style={
        rectangle, draw=black, fill=white, line width=0.8pt,
        minimum width=1.55cm, minimum height=0.52cm,
        font=\scriptsize, align=center, inner sep=0.8pt
    },
    flow/.style={-{Stealth[scale=0.55]}, line width=0.65pt},
    elabel/.style={font=\tiny, fill=white, inner sep=0.8pt, align=center},
    coltitle/.style={font=\scriptsize\bfseries, text=black!65}
]

%======================================================================
% DFD cấp 1 – 2 CỘT × 3 HÀNG (vừa A4 dọc, không chồng chéo)
% Mỗi khối: Thực thể → Tiến trình → Kho DL
%======================================================================

%========== CỘT TRÁI: 1.0 – 2.0 – 3.0 ==========
\node[coltitle] at (3.2, 9.6) {Nhóm chức năng cốt lõi};

%--- 1.0 ---
\node[entity]  (e1)  at (0.0, 8.5) {Quản trị viên};
\node[process] (p1)  at (3.5, 8.5) {1.0 Quản lý người dùng};
\node[store]   (d1)  at (6.6, 8.85) {D1 Users};
\node[store]   (d10) at (6.6, 8.15) {D10 Log};
\draw[flow] (e1) -- (p1) node[elabel, midway, above] {Cấu hình};
\draw[flow] (p1.east) -- ++(0.25,0) |- (d1.west);
\draw[flow] (p1.east) -- ++(0.25,0) |- (d10.west);

%--- 2.0 ---
\node[entity]  (e2) at (0.0, 6.5) {Nhân viên kho};
\node[process] (p2) at (3.5, 6.5) {2.0 Quản lý danh mục};
\node[store]   (d2) at (6.6, 6.85) {D2 HangHoa};
\node[store]   (d3) at (6.6, 6.15) {D3 NCC};
\draw[flow] (e2) -- (p2) node[elabel, midway, above] {QL danh mục};
\draw[flow] (p2.east) -- ++(0.25,0) |- (d2.west);
\draw[flow] (p2.east) -- ++(0.25,0) |- (d3.west);

%--- 3.0 ---
\node[entity]  (e3a) at (0.0, 4.7) {Nhân viên kho};
\node[entity]  (e3b) at (0.0, 3.9) {Quản lý kho};
\node[entity]  (e3c) at (0.0, 3.1) {Nhà cung cấp};
\node[process] (p3)  at (3.5, 3.9) {3.0 Quản lý nhập kho};
\node[store]   (d6)  at (6.6, 4.25) {D6 PNK};
\node[store]   (d8a) at (6.6, 3.55) {D8 TonKho};
\draw[flow] (e3a) -- (p3.west) node[elabel, pos=0.55, above] {Lập phiếu nhập};
\draw[flow] (e3b) -- (p3)        node[elabel, midway, above] {Duyệt};
\draw[flow] (e3c) -- (p3.west)   node[elabel, pos=0.55, below] {Giao hàng};
\draw[flow] (p3.east) -- ++(0.25,0) |- (d6.west);
\draw[flow] (p3.east) -- ++(0.25,0) |- (d8a.west)
    node[elabel, pos=0.7, below] {Tăng tồn};

%========== CỘT PHẢI: 4.0 – 5.0 – 6.0 ==========
\node[coltitle] at (12.0, 9.6) {Xuất kho / Tồn kho / Báo cáo};

%--- 4.0 ---
\node[entity]  (e4a) at (8.8, 8.5) {Nhân viên kho};
\node[entity]  (e4b) at (8.8, 7.7) {Quản lý kho};
\node[entity]  (e4c) at (8.8, 6.9) {Khách hàng};
\node[process] (p4)  at (12.3, 7.7) {4.0 Quản lý xuất kho};
\node[store]   (d7)  at (15.4, 8.15) {D7 PXK};
\node[store]   (d4)  at (15.4, 7.55) {D4 KH};
\node[store]   (d8b) at (15.4, 6.95) {D8 TonKho};
\draw[flow] (e4a) -- (p4.west) node[elabel, pos=0.55, above] {Lập phiếu xuất};
\draw[flow] (e4b) -- (p4)       node[elabel, midway, above] {Duyệt};
\draw[flow] (e4c) -- (p4.west)  node[elabel, pos=0.55, below] {Đơn hàng};
\draw[flow] (p4.east) -- ++(0.25,0) |- (d7.west);
\draw[flow] (p4.east) -- ++(0.25,0) |- (d4.west);
\draw[flow] (p4.east) -- ++(0.25,0) |- (d8b.west)
    node[elabel, pos=0.7, below] {Giảm tồn};

%--- 5.0 ---
\node[entity]  (e5) at (8.8, 5.1) {Nhân viên kho};
\node[process] (p5) at (12.3, 5.1) {5.0 Quản lý tồn kho};
\node[store]   (d9) at (15.4, 5.50) {D9 KiemKe};
\node[store]   (d5) at (15.4, 4.90) {D5 Kho};
\node[store]   (d8c) at (15.4, 4.30) {D8 TonKho};
\draw[flow] (e5) -- (p5) node[elabel, midway, above] {Kiểm kê};
\draw[flow] (p5.east) -- ++(0.25,0) |- (d9.west);
\draw[flow] (p5.east) -- ++(0.25,0) |- (d5.west);
\draw[flow] (p5.east) -- ++(0.25,0) |- (d8c.west);

%--- 6.0 ---
\node[entity]  (e6) at (8.8, 2.7) {Quản lý kho};
\node[process] (p6) at (12.3, 2.7) {6.0 Báo cáo thống kê};
\node[store]   (d6b) at (15.4, 3.15) {D6 PNK};
\node[store]   (d7b) at (15.4, 2.55) {D7 PXK};
\node[store]   (d8d) at (15.4, 1.95) {D8 TonKho};
\draw[flow] (e6) to[bend left=18] (p6)
    node[elabel, pos=0.55, above] {Yêu cầu BC};
\draw[flow] (p6) to[bend left=18] (e6)
    node[elabel, pos=0.55, below] {Báo cáo};
\draw[flow] (p6.east) -- ++(0.25,0) |- (d6b.west);
\draw[flow] (p6.east) -- ++(0.25,0) |- (d7b.west);
\draw[flow] (p6.east) -- ++(0.25,0) |- (d8d.west);

%----- Đường phân cột -----
\draw[gray!40, dashed, line width=0.6pt] (7.7, 9.3) -- (7.7, 1.6);

%----- Chú thích -----
\node[font=\scriptsize, anchor=west, align=left, text width=15.5cm] at (-0.5, 0.9) {%
  \textbf{Ký hiệu:}
  ellipse = thực thể ngoài;\quad
  HCN xám = tiến trình;\quad
  HCN trắng = kho dữ liệu (D).\\
  \textit{Ghi chú:} kho dữ liệu dùng chung (ví dụ D8) được lặp lại tại các tiến trình liên quan để tránh đường chéo.
};

\end{tikzpicture}%
}
\caption{DFD cấp 1 – Phân rã các luồng dữ liệu trong hệ thống quản lý kho hàng}
\label{fig:dfd-level1}
\end{figure}
